% ---- Abstract (<= half a page; excluded from the 3000-word limit) ----
\begin{abstract}
\noindent
Titan presents a rare combination of a dense, organic-rich atmosphere, stable surface liquids, an
abundant supply of prebiotic feedstock, and a probable subsurface water ocean, making it one of the
foremost astrobiological targets in the Solar System. Yet its habitability is constrained by the
physical separation of liquid water from surface organics, and by growing uncertainty over whether a
subsurface ocean survives at all. This group report synthesises five complementary studies that
together span Titan from its silicate core to its upper atmosphere: the tidal heating and thermal
evolution of the interior ocean; impact-driven melt production within the ice shell; prebiotic
reaction chemistry in transient impact brine pockets; longitudinal, seasonal structure in the
atmosphere; and a Bayesian, time-resolved map of surface habitability. Read together, these studies
identify impact cratering as a key process bridging Titan's isolated reservoirs, and geological time
as the axis along which its habitability rises and falls. We find that Titan offers multiple, largely
independent candidate habitats -- surface, ice-shell, and interior -- whose viability does not stand
or fall on the ocean question alone. We close with a combined outlook for Titan's habitability and the
tests that the forthcoming Dragonfly mission will provide.
\end{abstract}
